%% kapa: exam=1 round=22 year=2011 period=1 form=A subject=civil_law subject_ko=민법 seq=17 number=17
%% source: Q-Net 감정평가사 기출문제 (cst00309) artlSeq=5036790 file=extracted/1008/5036790/2011년 제22회 감정평가사 1차시험 기출문제/1교시.hwp
%% conversion: kordoc 4.24.51 -> tools/gen1.py (auto); review=proofread-sonnet; flags=none
\begin{question}{17}
乙은 甲으로부터 2011년 1월 5일 甲소유의 X토지를 매수하는 계약을 체결하면서 재산상황을 은폐하기 위하여 X토지에 대한 소유권이전등기를 여자친구 丙명의로 하기로 丙과 명의신탁약정을 하였다. 그 후 甲은 乙의 부탁대로 丙명의로 소유권이전등기를 해 주었다. 다음 설명으로 옳지 않은 것은? (다툼이 있으면 판례에 의함)
\begin{choices}
  \item 乙과 丙 사이의 약정은 무효이다.
  \item X토지에 대한 소유자는 여전히 甲이다.
  \item 乙은 丙에 대하여 명의신탁약정을 해지하고 X토지에 대한 소유권이전등기를 청구할 수 있다.
  \item 丙이 X토지를 丁에게 매도하고 소유권이전등기를 해 준 경우, 丁은 소유권을 취득한다.
  \item 乙이 丙과 혼인한 경우, 丙명의의 등기가 조세포탈, 강제집행의 면탈 또는 법령상 제한의 회피를 목적으로 한 것이 아니라면 혼인 후부터 그 등기는 유효하다.
\end{choices}
\end{question}
\answer{3}
